\section{Rubrik Komprehensif dan Penilaian}
\label{sec:ch17-rubrik}

\subsection{Dimensi Penilaian}

Rubrik komprehensif menilai lima dimensi:
\begin{enumerate}[label=\arabic*.,leftmargin=*]
  \item \textbf{pengetahuan}: ketepatan konsep dan istilah;
  \item \textbf{aplikasi}: kemampuan menerapkan teknik;
  \item \textbf{analisis}: kemampuan menganalisis masalah dan
        trade-off;
  \item \textbf{integrasi}: kemampuan menggabungkan komponen;
  \item \textbf{komunikasi}: kejelasan laporan dan presentasi.
\end{enumerate}

\begin{table}[H]
\centering
\caption{Rubrik komprehensif}
\label{tab:ch17-rubrik}
\begin{tabularx}{\textwidth}{@{}l X l@{}}
\toprule
\textbf{Dimensi} & \textbf{Indikator} & \textbf{Bobot} \\
\midrule
Pengetahuan & Konsep benar, istilah tepat & 10\% \\
Aplikasi & Teknik diterapkan dengan benar & 25\% \\
Analisis & Trade-off dan konflik dianalisis & 20\% \\
Integrasi & Komponen bekerja bersama & 30\% \\
Komunikasi & Laporan dan demo jelas & 15\% \\
\bottomrule
\end{tabularx}
\end{table}

\subsection{Level Kinerja}

Setiap dimensi dinilai pada empat level:
\begin{itemize}[leftmargin=*]
  \item \textbf{sangat baik}: lengkap, tepat, dan mendalam;
  \item \textbf{baik}: tepat dengan sedikit kekurangan;
  \item \textbf{cukup}: memenuhi kebutuhan dasar;
  \item \textbf{kurang}: ada kesalahan konsep atau bagian hilang.
\end{itemize}

\subsection{Konversi Nilai}

\begin{table}[H]
\centering
\caption{Konversi nilai akhir}
\label{tab:ch17-nilai}
\begin{tabular}{@{}ll@{}}
\toprule
\textbf{Rentang} & \textbf{Nilai} \\
\midrule
85--100 & A \\
70--84 & B \\
60--69 & C \\
50--59 & D \\
$<$50 & E \\
\bottomrule
\end{tabular}
\end{table}

\subsection{Menghindari Perhitungan Ganda}

Nilai praktikum dan proyek akhir harus dihitung secara terpisah.
Praktikum menilai milestone; proyek menilai integrasi akhir. Jika
mahasiswa menggunakan kode yang sama untuk keduanya, penilai harus
memastikan tidak ada pengulangan bobot.

\begin{warningbox}
Dokumentasikan pemetaan nilai dari setiap komponen ke CPMK. Tanpa
pemetaan yang jelas, sulit membuktikan ketercapaian kompetensi secara
akuntabel.
\end{warningbox}
